\VerdierPhysicalPageAnchor{0024}%
déduit pour les objets de cohomologie. Le mécanisme des suites spectrales cache
le plus souvent la simplicité des phénomènes. L'exemple élémentaire des
relations de Künneth est à cet égard significatif.

\VerdierSemanticAnchor{intro-5}
\noindent\textbf{5.}\quad
Soient $X$ et $Y$ deux espaces compacts, $A$ un anneau à élément unité
commutatif (pour simplifier), $F$ et $G$ deux faisceaux de $A$-modules sur $X$
et $Y$ respectivement. Lorsque $A$ est un corps, on sait que le produit
cartésien [23] définit un isomorphisme d'espaces vectoriels gradués :

\VerdierSemanticAnchor{intro-5-kunneth-field}
\begin{equation}
\mathrm H^*(X,F)\otimes_A\mathrm H^*(Y,G)
\xrightarrow{\sim}
\mathrm H^*(X\times Y,F\otimes_A G).
\tag{$\text{Kü}_{\mathrm{coh}}^{1}$}
\end{equation}

Lorsque $A=\mathbb Z$ et que le faisceau $F$ ou le faisceau $G$ est sans
torsion, on a des suites exactes scindables :

\VerdierSemanticAnchor{intro-5-kunneth-torsion-free}
\begin{equation}
\begin{aligned}
0 &\longrightarrow
\bigoplus_{p+q=n}\mathrm H^p(X,F)\otimes_A\mathrm H^q(Y,G)
\longrightarrow \mathrm H^n(X\times Y,F\otimes_A G) \longrightarrow \\
&\longrightarrow
\bigoplus_{p+q=n+1}
\operatorname{Tor}^{\mathbb Z}\bigl(\mathrm H^p(X,F),\mathrm H^q(Y,G)\bigr)
\longrightarrow 0.
\end{aligned}
\tag{$\text{Kü}_{\mathrm{coh}}^{2}$}
\end{equation}

Lorsque $A=\mathbb Z$ et qu'on ne fait aucune hypothèse sur $F$ et $G$, on a
deux familles de suites exactes scindables :

\VerdierSemanticAnchor{intro-5-kunneth-general-split}
\begin{equation}
\begin{aligned}
0 &\longrightarrow
\bigoplus_{p+q=n}\mathrm H^p(X,F)\otimes_A\mathrm H^q(Y,G)
\longrightarrow \mathbb H^n \longrightarrow \\
&\longrightarrow
\bigoplus_{p+q=n+1}
\operatorname{Tor}^{\mathbb Z}\bigl(\mathrm H^p(X,F),\mathrm H^q(Y,G)\bigr)
\longrightarrow 0, \\
0 &\longrightarrow
\mathrm H^{n+1}\bigl(X\times Y,\mathcal T\!or^A(F,G)\bigr)
\longrightarrow \mathbb H^n
\longrightarrow \mathrm H^n(X\times Y,F\otimes_A G)
\longrightarrow 0.
\end{aligned}
\tag{$\text{Kü}_{\mathrm{coh}}^{3}$}
\end{equation}

La formulation initiale $\text{(Kü)}_{\mathrm{coh}}^{1}$ se
complique donc, mais reste accessible au calcul explicite. Quand l'anneau $A$
est plus compliqué, par exemple lorsque
$A=\mathbb Z/\ell^r\mathbb Z$, $\ell$ premier, $r>1$, anneau de dimension
cohomologique globale infinie, on n'a plus, traditionnellement, que deux
suites spectrales ayant même aboutissement dont les termes initiaux sont :

\VerdierSemanticAnchor{intro-5-kunneth-spectral-sequences}
\begin{equation}
\begin{aligned}
{}'K_2^{p,q}
  &=\bigoplus_{r+s=q}
    \operatorname{Tor}_{-p}^{A}
    \bigl(\mathrm H^r(X,F),\mathrm H^s(Y,G)\bigr), \\
{}''K_2^{p,q}
  &=\mathrm H^p\bigl(X\times Y,\mathcal T\!or_{-q}^{A}(F,G)\bigr).
\end{aligned}
\tag{$\text{Kü}_{\mathrm{coh}}^{4}$}
\end{equation}

Cependant ces suites spectrales et l'isomorphie entre leurs aboutissements ne
sont que la conséquence et la traduction imparfaite de la relation suivante
dans la catégorie dérivée de la catégorie des $A$-modules :

\VerdierSemanticAnchor{intro-5-kunneth-derived}
\begin{equation}
\mathrm{RH}(X,F)\overset{\mathrm L}{\otimes}_A\mathrm{RH}(Y,G)
\xrightarrow{\sim}
\mathrm{RH}\bigl(X\times Y,F\overset{\mathrm L}{\otimes}_A G\bigr),
\tag{$\text{Kü}_{\mathrm{Der}}$}
\end{equation}
